\chapter{Literature Review}
\label{chap:lr}
\chaptermark{Literature Review}

Several literature sources were employed to perform analysis of the problem. Google Scholar and Research Gate served to perform the research on the distributed compute management and pull-model approach. The keywords used to find the sources were "cloud orchestration", "distributed cloud computing", "pull-model". To fetch more practical insights we decided to use official Kubernetes, HashiCorp Nomad and Apache OpenServerless documentation.

\section{Container orchestration systems}

Modern orchestration platforms, such as Kubernetes and HashiCorp Nomad, have become standard tools for managing distributed workloads in cloud environments. Kubernetes provides a declarative configuration model and uses a centralized control plane to manage containerized workloads across multiple nodes \cite{k8s_docs}. However, it relies on hybrid (both push and pull models) control mechanism. To decide where to schedule workloads, there is a push-based mechanism, located on the control plane, which could serve as another point of failure. But for configuration delivery it uses pull model, node agent watches for changes and applies it when there are some. Comprehensive surveys of Kubernetes scheduling \cite{carrion2022k8ssurvey, senjab2023k8salgorithms} classify its scheduling strategies into filter-score-prioritize pipelines and identify open challenges in multi-cluster and edge deployments. Nomad offers a similar architecture but emphasizes simplicity and flexibility in scheduling non-containerized workloads \cite{nomad_docs}. Nevertheless agent or client in terms of Nomad acts like a reporter of a node current state whether a resource operator. In conclusion both rely primarily on a push-based control mechanism, where the control plane assigns tasks to worker nodes. While this approach ensures strong control and consistency, it introduces higher coupling, operational complexity, and reduced fault tolerance.

Similarly, research on service mesh architectures \cite{canal2024mesh} demonstrates that eliminating sidecar proxies through novel control plane designs can significantly reduce control plane overhead compared to traditional approaches like Istio. This sidecar-free architecture for multi-tenant environments provides a viable alternative to heavily-layered orchestration systems, aligning with the lightweight design principles explored in this thesis.

Research by Sebrechts et al. \cite{sebrechts2018orchestrator} highlights challenges in distributed application management, emphasizing the need for scalable and adaptive orchestration systems. It offers great alternative - distributing operational agents. Similarly, Pantazoglou et al. \cite{pantazoglou2015decentralized} propose decentralized workload management approaches to improve energy efficiency and reduce coordination overhead in large-scale systems. These studies suggest that centralization, while powerful, can become a bottleneck for resilience and scalability.

The evolution of cluster scheduling architectures further illustrates this tension. Google's Borg \cite{verma2015borg} employs a centralized monolithic scheduler with per-machine agents (Borglets), while its successor Omega \cite{schwarzkopf2013omega} introduced shared-state optimistic concurrency to improve scalability. Apache Mesos \cite{hindman2011mesos} and YARN \cite{vavilapalli2013yarn} adopted two-level scheduling, decoupling resource management from application-specific scheduling logic. More recently, Sparrow \cite{ousterhout2013sparrow} demonstrated that stateless, distributed scheduling via random sampling can achieve low-latency placement for short-lived tasks, and Apollo \cite{boutin2015apollo} explored coordinated decentralized scheduling with per-job managers. Hawk \cite{delgado2015hawk} proposed a hybrid approach, delegating short tasks to distributed schedulers while retaining centralized control for long-running jobs. Firmament \cite{gog2016firmament} applied min-cost max-flow optimization to centralized cluster scheduling. These systems collectively demonstrate a progression from monolithic to shared-state to fully distributed designs, yet all retain some form of centralized coordination that the pull-based model proposed in this thesis seeks to minimize.

\section{Pull-based and decentralized approaches}

The pull-based paradigm, originally explored in software engineering and version control systems, introduces an alternative model for coordination and synchronization. Recent work on pull-based scheduling for serverless computing \cite{hiku2025pull} demonstrates that eliminating centralized scheduling decisions in favor of worker-initiated workload acquisition can reduce coordination overhead and improve scalability. These principles can be adapted to workload orchestration, where worker nodes autonomously retrieve workloads and configurations, thus minimizing reliance on a central scheduler.

Recent research on decentralized cloud computing supports this shift. Pantazoglou et al. \cite{pantazoglou2015decentralized} propose decentralized and energy-efficient workload management approaches that reduce coordination overhead in enterprise clouds, while Hiku \cite{hiku2025pull} demonstrates the viability of pull-based scheduling in serverless environments. These approaches align with modern distributed architectures that prioritize resilience and fault isolation over strict central coordination.

The CAP theorem \cite{brewer2012cap, gilbert2012cap} formalizes the fundamental trade-off between consistency, availability, and partition tolerance in distributed systems. Brewer \cite{brewer2012cap} argues that modern systems can mitigate CAP limitations through techniques such as eventual consistency and conflict-free replicated data types, which aligns with the design philosophy of the proposed system. Amazon's Dynamo \cite{decandia2007dynamo} exemplifies an availability-first, eventually consistent design that inspired many cloud-native architectures. In contrast, traditional coordination services such as ZooKeeper \cite{hunt2010zookeeper} and Google's Chubby \cite{burrows2006chubby} provide strong consistency through consensus-based distributed locking, while Paxos \cite{lamport2001paxos} and Raft \cite{ongaro2014raft} offer theoretically grounded consensus protocols at the cost of increased coordination overhead. The pull-based orchestration model proposed in this thesis deliberately avoids such consensus requirements: by delegating scheduling decisions to autonomous agents that poll for available workloads, the system favors availability and partition tolerance over strong consistency, a trade-off justified by the workload characteristics and operational requirements identified in the problem analysis.

Comprehensive surveys of geo-distributed task scheduling \cite{wu2025task} further emphasize that distributed systems operating across heterogeneous network conditions and region-specific resource constraints require sophisticated coordination approaches. Such systems must balance performance enhancement, fairness assurance, and fault tolerance improvement while managing challenges inherent to geographic distribution, including varying network latency and computational capabilities across locations.


\section{Security and reliability concerns}

Infrastructure reliability and security are critical for large-scale deployments. Google Cloud’s production systems rely on strong security models and automated workload verification to ensure service integrity \cite{google_security}. However, such designs also depend on highly centralized control planes, which can create single points of failure. Research on resilient cloud systems emphasizes the need for autonomous components that can recover independently \cite{liaqat2017federated} and sustain partial system failures \cite{velev2011cloud, rasheed2014data}. Finally, compliance and governance are integral to reliable infrastructure, especially in enterprise and regulated environments. Emerging research also highlights the potential of decentralized and federated cloud management frameworks to enhance resilience and security while reducing reliance on single points of failure.

\section{Lightweight and edge orchestration}

The demand for lightweight orchestration has intensified with the proliferation of edge computing and resource-constrained IoT environments. Čilić et al. \cite{cilic2023edgeeval} benchmark Kubernetes, K3s, and KubeEdge on edge hardware, demonstrating that even lightweight Kubernetes distributions incur significant overhead relative to bare-metal execution. Usman et al. \cite{usman2024lightweight} extend this analysis to lightweight distributions including K3s and k0s, finding that resource consumption remains substantial for small-scale edge deployments. Böhm and Wirtz \cite{bohm2022cloudedge} survey Kubernetes-based orchestration architectures for smart city edge scenarios and identify architectural complexity as a barrier to adoption in resource-constrained settings.

Alternative frameworks have emerged to address these limitations. FogBus2 \cite{deng2021fogbus2} provides a lightweight container-based framework integrating IoT, edge, and cloud resources with minimal operational overhead. Arora et al. \cite{arora2022leso} propose a lightweight edge slice orchestration framework for multi-access edge computing with a minimal control plane. Gand et al. \cite{gand2020fuzzy} introduce a self-adaptive fuzzy controller for lightweight edge container orchestration, demonstrating that autonomous reconciliation mechanisms can maintain service quality without centralized scheduling. Wang et al. \cite{wang2022edgefog} examine container orchestration in edge and fog environments for real-time IoT applications, emphasizing the need for low-overhead scheduling mechanisms. Morabito et al. \cite{morabito2025consensus} explore consensus-based distributed orchestration for edge microservices, though their approach introduces coordination overhead that the pull-based model seeks to avoid.

\section{Fault tolerance in container-based systems}

Fault tolerance mechanisms in container orchestration span failure detection, state recovery, and workload rescheduling. Bakhshi and Rodriguez-Navas \cite{bakhshi2021faulttolerant} study fault-tolerant permanent storage for container-based fog architectures, highlighting that stateless compute layers combined with durable storage backends provide an effective resilience pattern. The seminal work on leases by Gray and Cheriton \cite{gray1989leases} establishes time-based distributed locking as a lightweight alternative to consensus-based coordination, a principle directly applicable to workload claiming in pull-based systems. MapReduce \cite{dean2004mapreduce} introduced the pattern of automatic task re-execution upon failure, inspiring the reconciliation-based fault tolerance mechanisms adopted in modern orchestrators. TetriSched \cite{grandl2016tetrisched} addresses adaptive rescheduling in shared heterogeneous clusters, demonstrating that dynamic plan revision can improve placement quality, albeit at the cost of increased scheduling complexity.

\section{Summary}

The reviewed literature provides valuable information on modern approaches to workload management, orchestration and security in cloud environments. Existing systems such as Kubernetes and Nomad demonstrate effective solutions but still rely heavily on centralized management models rooted in decades of cluster scheduling research \cite{verma2015borg, schwarzkopf2013omega, hindman2011mesos, vavilapalli2013yarn}. Research shows that decentralized and pull-based mechanisms can increase scalability, resilience and autonomy. CAP theory \cite{brewer2012cap, gilbert2012cap} and the design of availability-first systems \cite{decandia2007dynamo} provide theoretical justification for architectures that favor partition tolerance over strong consistency. Studies on lightweight edge orchestration \cite{cilic2023edgeeval, usman2024lightweight, deng2021fogbus2, arora2022leso} demonstrate demand for minimal-footprint scheduling solutions. Overall, while much can be gained from the work already done, there is still considerable room for improvement - especially in developing more adaptive, fault-tolerant and secure orchestration architectures.
